\section{Discussion and Conclusion}
\label{sec:conclusion}

\textsc{MWG-EW} is a bridge between neural architecture and runtime design:
the neural component makes a low-rank conditional descriptor useful on some
inputs; the systems component ensures the descriptor does not quietly become
another persistent tensor. The parameter-lifetime view---asking which
weight-like objects must cross durable memory and synchronization
boundaries---is the broader contribution.

\paragraph{What is established.}
Four claims are measured (Section~\ref{sec:results}): the fused no-store
arm at latency parity with materialization on Ascend~910B2 with no core
write-back; same-rank FFN-layer reconstruction parity of the generated
diagonal core
(r384, r401, r512, four seeds, one coherent run family), with a
$+1.95\%$ equal-budget penalty (persistent r401 vs.\ causal r384 at
48.2 vs.\ 48.1~MiB) that same-rank parity does not close; recompute-in-backward cutting the training activation
footprint $24.5\times$--$71.8\times$ at exact gradient equality; and the
optimizer-inclusive synchronized training step running
$3.37\times$--$4.14\times$ faster than dense at 8B-scale geometry---a
low-rank parameter-count effect on the training side, distinct from the
inference-scoped no-store contract.

\paragraph{What remains open.}
Three validation gaps separate this prototype from a finished systems
deployment (enumerated in Section~\ref{sec:results}); each could
invalidate the core bandwidth-reduction thesis if it does not resolve
favorably: the operational router is unrealized (a
\texttt{torch.where} fallback evaluates both branches, and without
kernel-level conditional execution the benefit collapses to the
unconditional low-rank case); the beyond-L2 epilogue bottleneck is now
resolved at the kernel level (the measured v4 fix of
Section~\ref{sec:systems} is bit-identical to v1 and up to $3.03\times$
faster at $T{=}4096$, clearing the ${\sim}2\times$ viability threshold),
but its deployment benefit is bounded by the unmeasured serving stack;
and deployment-relevant serving
regimes are unvalidated, so the net off-chip bandwidth reduction at
deployment scale remains an analytic projection rather than a measured
end-to-end serving win.

\paragraph{Why the negative results matter.}
The integrity audit that withdrew prior quality and routing claims is
part of the result: ephemerality is easy to conflate with ordinary low
rank, and a \texttt{torch.where} fallback can mask the absence of true
branch skipping.  The shuffled-context (Table~\ref{tab:shuffled}) and
rank-512 (Table~\ref{tab:rank512control}) controls add two
negative-leaning findings: context-conditioning contributes only
${\sim}0.05\%$ of MSE (paired $-0.00021$ at r384; not detectable at
r512), and a persistent r512 factor
beats the generated r384 core by $14.6\%$ relative MSE---rank is the
dominant quality lever, so the ephemeral core must justify itself on
bandwidth and lifetime grounds, not quality grounds.  Future work should
pre-register disjoint evaluations, distinguish token-level patching from
whole-call dispatch, and measure physical traffic rather than logical
payload counts.
